% Generated consolidated text; do not edit manually.
% Source records are preserved in the project archive.

\subsection*{Convocatoria para la Atención de Proyectos Estratégicos de Generación y Almacenamiento de Energía Eléctrica}


Versión consolidada completa de consulta, integrada con la Convocatoria original publicada en el Diario Oficial de la Federación el 15 de mayo de 2026, la modificación publicada el 26 de mayo de 2026 y la segunda modificación publicada el 10 de julio de 2026.\referenciaBase{}\referenciaPrimera{}\referenciaSegunda{}


\subsubsection*{Nota editorial}


Este documento reproduce la estructura completa del instrumento para que pueda leerse de corrido como si fuera un solo texto, incorporando las modificaciones en su lugar correspondiente y señalando cada ajuste con marcas editoriales.\referenciaBase{}\referenciaPrimera{}\referenciaSegunda{}


No sustituye las publicaciones oficiales del Diario Oficial de la Federación, sino que funciona como versión consolidada de consulta para revisión, trazabilidad normativa y edición posterior en otros formatos, incluido PDF.\referenciaBase{}\referenciaPrimera{}\referenciaSegunda{}


\subsubsection*{Convención de marcas}


\begin{longtable}{@{}>{\raggedright\arraybackslash}p{0.25\textwidth}>{\raggedright\arraybackslash}p{0.735\textwidth}@{}}
\toprule
\senerth{Marca} & \senerth{Significado} \\
\midrule
\endfirsthead
\toprule
\senerth{Marca} & \senerth{Significado} \\
\midrule
\endhead
\marcaPrimera{} & Texto modificado por Acuerdo publicado en el DOF el 26/05/2026.\referenciaPrimera{} \\
\marcaSegunda{} & Texto modificado por segunda modificación publicada en el DOF el 10/07/2026.\referenciaSegunda{} \\
\marcaTercera{} & Texto adicionado por segunda modificación publicada en el DOF el 10/07/2026.\referenciaSegunda{} \\
Sin marca & Texto que conserva la redacción de la Convocatoria original publicada el 15/05/2026.\referenciaBase{} \\
\bottomrule
\end{longtable}


\subsubsection*{Encabezado}


CONVOCATORIA para la Atención de Proyectos Estratégicos de Generación y Almacenamiento de Energía Eléctrica, alineados a la planeación vinculante.\referenciaBase{}


\subsubsection*{Fundamento y considerandos integrados}


La Convocatoria original fue emitida por la Secretaría de Energía con fundamento en disposiciones constitucionales, orgánicas, administrativas y sectoriales relacionadas con la rectoría del Estado en el sector eléctrico, la planeación vinculante, el otorgamiento de permisos de generación y almacenamiento, y los esquemas de participación de particulares y de asociación con la CFE.\referenciaBase{}


En sus considerandos, la Convocatoria explica el contexto de reforma constitucional y legal del sector eléctrico, el fortalecimiento del papel de la Nación en la planeación y control del Sistema Eléctrico Nacional, el nuevo arreglo de atribuciones de la SENER y la CNE, y la necesidad de incorporar generación renovable y almacenamiento para fortalecer la confiabilidad, continuidad, eficiencia, seguridad y sostenibilidad del SEN.\referenciaBase{}


La primera modificación justificó el ajuste de fechas del calendario por el interés manifestado por las personas participantes y por la conveniencia de favorecer una mayor participación en la Convocatoria.\referenciaPrimera{}


La segunda modificación explicó que la revisión integral del procedimiento buscó incentivar la participación, precisar etapas, fortalecer la revisión por autoridades y entes participantes, dar mayor claridad al proceso y evitar costos innecesarios a los particulares, sin incorporar requisitos adicionales de elegibilidad ni modificar el tipo de proyectos susceptibles de participar.\referenciaSegunda{}


\subsubsection*{Contenido}


\puntoConvocatoria{1. Objetivos y alcance.\referenciaBase{}}

\puntoConvocatoria{2. Glosario.\referenciaBase{}}

\puntoConvocatoria{3. Reglas generales y de interpretación.\referenciaBase{}}

\puntoConvocatoria{4. Ventanilla Única de Proyectos Estratégicos del Sector Energético.\referenciaBase{}}

\puntoConvocatoria{5. Participación de las distintas autoridades en la Convocatoria de Proyectos Estratégicos y consideraciones sobre sus trámites.\referenciaBase{}}

\puntoConvocatoria{6. Etapas y Actividades.\referenciaBase{}}

\puntoConvocatoria{7. Estudios de conexión e interconexión y pagos correspondientes.\referenciaBase{}}

\puntoConvocatoria{8. Requisito de aceptación de Obras de Interconexión, Obras de Conexión y Obras de Refuerzo.\referenciaBase{}}

\puntoConvocatoria{9. Alineación a la planeación vinculante.\referenciaBase{}}

\puntoConvocatoria{10. Evaluación de los proyectos por parte del Grupo de Análisis Técnico.\referenciaBase{}}

\puntoConvocatoria{11. Procedencia de Proyectos en Asociación con la CFE.\referenciaBase{}}

\puntoConvocatoria{12. Aprobación de permisos y emisión de los títulos de permisos.\referenciaBase{}}

\puntoConvocatoria{13. Presentación y devolución de las garantías financieras asociadas a contratos de conexión e interconexión derivados de los permisos otorgados bajo la Convocatoria de Proyectos Estratégicos.\referenciaBase{}}

\puntoConvocatoria{14. Publicación periódica de capacidades y resultados de la Convocatoria de Proyectos Estratégicos.\referenciaBase{}}

\puntoConvocatoria{15. Suscripción del contrato de conexión o interconexión correspondiente a los permisos otorgados bajo la Convocatoria de Proyectos Estratégicos.\referenciaBase{}}

\puntoConvocatoria{16. Seguimiento de permisos hasta la Entrada en Operación Comercial de los proyectos.\referenciaBase{}}

\puntoConvocatoria{17. Etapas, plazos y actividades de la Convocatoria de Proyectos Estratégicos.\referenciaBase{}}

\puntoConvocatoria{18. Referencias de Capacidad Requerida de Sistemas de Almacenamiento de Energía Eléctrica y Generación de Energía Eléctrica.\referenciaBase{}}

\puntoConvocatoria{19. Anexos.\referenciaBase{}}


\subsubsection*{1. Objetivos y alcance}


\paragraph*{1.1 Objetivos}
\mbox{}\par
\nopagebreak[4]


La Convocatoria de Proyectos Estratégicos tiene los siguientes objetivos:\referenciaBase{}


\puntoConvocatoria{I. Establecer el mecanismo para la atención prioritaria de las solicitudes para la obtención de permisos, los estudios de interconexión y conexión y, en su caso, la celebración de los contratos de interconexión y conexión, que permita el desarrollo de proyectos estratégicos que se requieran para fortalecer la confiabilidad, continuidad, eficiencia, seguridad y sostenibilidad del SEN asociados a: a) Proyectos de Sistemas de Almacenamiento de Energía Eléctrica que se desarrollen bajo Proyecto en Asociación con la CFE, o b) Proyectos de generación de energía eléctrica renovable que se desarrollen bajo Proyecto en Asociación con la CFE.\referenciaBase{}}


\puntoConvocatoria{II. Fomentar el desarrollo de nuevos proyectos estratégicos alineados con la planeación vinculante establecida por la SENER en atención a las necesidades energéticas del país.\referenciaBase{}}


\puntoConvocatoria{III. Ordenar y agilizar los trámites de registro, evaluación y seguimiento de los proyectos asociados a las solicitudes de permisos que se alineen con la planeación vinculante y se desarrollen conforme a esta Convocatoria de Proyectos Estratégicos, a través de sus distintas etapas, entre las que se encuentran la evaluación y aprobación de las solicitudes por parte de la Comisión Nacional de Energía.\referenciaBase{}}


En ese sentido, esta Convocatoria de Proyectos Estratégicos no establece requisitos adicionales a los previstos en la normatividad aplicable, sino que define un conjunto mínimo de elementos que permiten diferenciar las solicitudes con base en su solidez técnica, su alineación con la planeación vinculante y su viabilidad técnica y operativa.\referenciaBase{}


\puntoConvocatoria{IV. Establecer la Capacidad Requerida para la Convocatoria de Proyectos Estratégicos para los proyectos que se presenten conforme a los términos de la Convocatoria de Proyectos Estratégicos, considerando las necesidades energéticas del Sistema Eléctrico Nacional de acuerdo con la planeación vinculante.\marcaSegunda{}\referenciaSegunda{}}


\puntoConvocatoria{V. Establecer el procedimiento para que las personas solicitantes registren sus solicitudes bajo Proyectos en Asociación con la CFE.\marcaSegunda{}\referenciaSegunda{}}


\puntoConvocatoria{VI. Establecer que la CNE es la responsable de recibir, tramitar y resolver las solicitudes de permisos a través de su OPE, para lo cual debe dar seguimiento a la emisión de los resultados de los estudios de conexión e interconexión por parte del CENACE.\referenciaBase{}}


\paragraph*{1.2 Alcance}
\mbox{}\par
\nopagebreak[4]


La Convocatoria de Proyectos Estratégicos está dirigida a todas las personas solicitantes interesadas en desarrollar en los términos previstos en la misma los siguientes proyectos:\referenciaBase{}


\puntoConvocatoria{I. Centrales de generación de energía eléctrica renovable con capacidad igual o mayor a 0.7 MW, o\referenciaBase{}}

\puntoConvocatoria{II. Sistemas de Almacenamiento de Energía Eléctrica no Asociados con capacidad igual o mayor a 0.7 MW.\referenciaBase{}}


\paragraph*{1.2.1}
\mbox{}\par
\nopagebreak[4]


Los proyectos y solicitudes que sean presentados por las personas solicitantes en los términos de esta Convocatoria deben cumplir con las siguientes características:\referenciaBase{}


\puntoConvocatoria{I. Contribuir al cumplimiento de los objetivos de confiabilidad, eficiencia y sostenibilidad del SEN.\referenciaBase{}}

\puntoConvocatoria{II. Contribuir a las necesidades energéticas del país conforme a lo establecido en el numeral 18 de esta Convocatoria de Proyectos Estratégicos.\referenciaBase{}}

\puntoConvocatoria{III. Contar con la evidencia documental de los trámites de viabilidad técnica, social y ambiental efectuados, en su caso, ante las autoridades competentes, tales como la SENER, el CENACE y la Secretaría de Medio Ambiente y Recursos Naturales.\referenciaBase{}}

\puntoConvocatoria{IV. Manifestar el Esquema de Desarrollo Mixto o asociación bajo la cual pretende desarrollar el Proyecto en Asociación con la CFE.\referenciaBase{}}


\paragraph*{1.2.2}
\mbox{}\par
\nopagebreak[4]


Quedan excluidos de participar en la Convocatoria de Proyectos Estratégicos los siguientes proyectos:\referenciaBase{}


\puntoConvocatoria{I. Para el caso de generación de energía eléctrica:\referenciaBase{}}

a) Generación distribuida.\referenciaBase{}

b) Autoconsumo en sus diferentes modalidades.\referenciaBase{}

c) Generación de energía eléctrica en modalidad de cogeneración.\referenciaBase{}

d) Centrales Eléctricas que, de forma previa a la publicación de esta Convocatoria, hayan obtenido un permiso de generación de energía eléctrica.\referenciaBase{}


\puntoConvocatoria{II. Para el caso de los Sistemas de Almacenamiento de Energía Eléctrica:\referenciaBase{}}

a) Sistemas de Almacenamiento de Energía Eléctrica que no requieran permiso o autorización conforme a la Ley del Sector Eléctrico y su reglamento.\referenciaBase{}

b) Sistemas de Almacenamiento de Energía Eléctrica asociados a la infraestructura destinada a la prestación del Servicio Público de Transmisión y Distribución de Energía Eléctrica.\referenciaBase{}

c) Sistemas de Almacenamiento de Energía Eléctrica asociados a Centrales Eléctricas o a Centros de Carga vinculados a permisos o autorizaciones vigentes a la fecha de presentación del registro de manifestación de interés correspondiente.\referenciaBase{}


\subsubsection*{2. Glosario}


\puntoConvocatoria{I. CENACE: Centro Nacional de Control de Energía.\referenciaBase{}}

\puntoConvocatoria{II. CNE: Comisión Nacional de Energía.\referenciaBase{}}

\puntoConvocatoria{III. CFE: Empresa Pública del Estado, Comisión Federal de Electricidad.\referenciaBase{}}

\puntoConvocatoria{IV. Capacidad de Almacenamiento: Cantidad de energía, expresada en MWh, que puede ser almacenada para ser entregada como energía eléctrica y mediante la cual se identifica al Sistema de Almacenamiento de Energía Eléctrica, indicada en corriente alterna.\referenciaBase{}}

\puntoConvocatoria{V. Capacidad Requerida para la Convocatoria de Proyectos Estratégicos: Para los efectos de la presente Convocatoria de Proyectos Estratégicos, la capacidad requerida, expresada en MW, comprende la capacidad de generación de energía eléctrica por sistema interconectado, así como la potencia de los SAEE por Gerencia de Control Regional. Dicha capacidad constituye el objetivo a cubrir mediante los proyectos que participen en la Convocatoria, a fin de atender las necesidades prioritarias del Sistema Eléctrico Nacional, y se cubrirá de manera progresiva conforme a lo previsto en la presente Convocatoria.\marcaSegunda{}\referenciaSegunda{}}

\puntoConvocatoria{VI. Convocatoria de Proyectos Estratégicos: La Convocatoria para la atención de Proyectos Estratégicos de generación y almacenamiento de energía eléctrica, alineados a la planeación vinculante.\marcaSegunda{}\referenciaSegunda{}}

\puntoConvocatoria{VII. Disposiciones de la Planeación Vinculante: Disposiciones administrativas de carácter general para la planeación vinculante aplicables a la generación de energía eléctrica publicadas en el Diario Oficial de la Federación el 17 de octubre de 2025 y, las Disposiciones administrativas de carácter general para la planeación vinculante aplicables a Almacenamiento de Energía Eléctrica que para tal efecto emita la SENER.\marcaSegunda{}\referenciaSegunda{}}

\puntoConvocatoria{VIII. Disposiciones de Permisos: Disposiciones administrativas de carácter general que establecen los términos legales, técnicos y financieros para solicitar el otorgamiento y la modificación de permisos de generación y de almacenamiento de energía eléctrica, así como su vigencia, publicadas en el Diario Oficial de la Federación el 23 de octubre de 2025.\marcaSegunda{}\referenciaSegunda{}}

\puntoConvocatoria{IX. Disposiciones de Integración de SAEE: Disposiciones Administrativas de Carácter General para la integración de Sistemas de Almacenamiento de Energía Eléctrica al Sistema Eléctrico Nacional publicadas en el Diario Oficial de la Federación el 16 de abril de 2026.\marcaSegunda{}\referenciaSegunda{}}

\puntoConvocatoria{X. Entrada en Operación Comercial: Fecha determinada en la planeación vinculante en la que un proyecto debe iniciar operaciones comerciales y que constituye el parámetro temporal obligatorio para validar criterios de firmeza para el desarrollo de los proyectos presentados en las solicitudes correspondientes.\marcaSegunda{}\referenciaSegunda{}}

\puntoConvocatoria{XI. Esquema de Desarrollo Mixto: Modalidad de participación para el desarrollo de proyectos entre la CFE y particulares para el desarrollo de proyectos de generación conforme a la Ley de la Empresa Pública del Estado, Comisión Federal de Electricidad, los lineamientos aplicables y los instrumentos que se suscriban para cada proyecto.\marcaSegunda{}\referenciaSegunda{}}

\puntoConvocatoria{XII. Grupo de Análisis Técnico o GAT: Grupo integrado en términos de las Disposiciones de la Planeación Vinculante.\marcaSegunda{}\referenciaSegunda{}}

\puntoConvocatoria{XIII. LCNE: Ley de la Comisión Nacional de Energía.\marcaSegunda{}\referenciaSegunda{}}

\puntoConvocatoria{XIV. LPTE: Ley de Planeación y Transición Energética.\marcaSegunda{}\referenciaSegunda{}}

\puntoConvocatoria{XV. LSE: Ley del Sector Eléctrico.\marcaSegunda{}\referenciaSegunda{}}

\puntoConvocatoria{XVI. MW: Megawatt.\marcaSegunda{}\referenciaSegunda{}}

\puntoConvocatoria{XVII. MIC: Manual de Interconexión de Centrales Eléctricas y Conexión de Centros de Carga publicado en el Diario Oficial de la Federación el 9 de febrero de 2018 o aquel que lo sustituya.\marcaSegunda{}\referenciaSegunda{}}

\puntoConvocatoria{XVIII. Nivel de Tensión: Magnitud de voltaje en el que la Central Eléctrica o el Sistema de Almacenamiento de Energía Eléctrica se conecta o interconecta, según corresponda, al SEN, expresado en kilovoltios.\marcaSegunda{}\referenciaSegunda{}}

\puntoConvocatoria{XIX. OPE: Oficialía de Partes Electrónica, herramienta tecnológica con la que cuenta la CNE, para recibir y dar seguimiento a las solicitudes de permisos presentadas en términos de la Convocatoria de Proyectos Estratégicos, así como para dar seguimiento a los permisos de generación y almacenamiento de energía eléctrica. Dicha herramienta se regula con base en lo establecido en las reglas generales para su funcionamiento publicadas en el Diario Oficial de la Federación, y se encuentra disponible en el portal \href{https://ope.cne.gob.mx/}{\texttt{enlace}}.\marcaSegunda{}\referenciaSegunda{}}

\puntoConvocatoria{XX. Potencia SAEE: Potencia activa, expresada en MW, e indicada en corriente alterna, mediante la cual se designa e identifica al Sistema de Almacenamiento de Energía Eléctrica.\marcaSegunda{}\referenciaSegunda{}}

\puntoConvocatoria{XXI. Punto de Interconexión: Punto físico de la Red Nacional de Transmisión o de las Redes Generales de Distribución donde una Central Eléctrica se interconecta físicamente al SEN para inyectar energía eléctrica o donde un Sistema de Almacenamiento de Energía Eléctrica se conecta al SEN para cargar energía eléctrica.\marcaSegunda{}\referenciaSegunda{}}

\puntoConvocatoria{XXII. Proyecto en Asociación con la CFE: Proyecto cuya ejecución se plantea con participación de la Comisión Federal de Electricidad, mediante un Esquema para el Desarrollo Mixto, sea Inversión Mixta, Producción de Largo Plazo, asociación, alianza o cualquier otra modalidad permitida por la Ley de la Empresa Pública del Estado, Comisión Federal de Electricidad, los Lineamientos de los Esquemas para el Desarrollo Mixto de la CFE y demás disposiciones aplicables.\marcaSegunda{}\referenciaSegunda{}}

\puntoConvocatoria{XXIII. RLPTE: Reglamento de la Ley de Planeación y Transición Energética.\marcaSegunda{}\referenciaSegunda{}}

\puntoConvocatoria{XXIV. RLSE: Reglamento de la Ley del Sector Eléctrico.\marcaSegunda{}\referenciaSegunda{}}

\puntoConvocatoria{XXV. SENER: Secretaría de Energía.\marcaSegunda{}\referenciaSegunda{}}

\puntoConvocatoria{XXVI. SEN: Sistema Eléctrico Nacional.\marcaSegunda{}\referenciaSegunda{}}

\puntoConvocatoria{XXVII. SIASIC: Sistema de Atención a Solicitudes de Interconexión y Conexión del CENACE.\marcaSegunda{}\referenciaSegunda{}}

\puntoConvocatoria{XXVIII. Subestación: Instalación de la RNT o de las RGD en donde un elemento de generación de energía eléctrica o Sistema de Almacenamiento de Energía Eléctrica se conecta o interconecta físicamente, según su condición operativa, al SEN, para cargar o inyectar energía eléctrica.\marcaSegunda{}\referenciaSegunda{}}

\puntoConvocatoria{XXIX. VUPE: Ventanilla Única de Proyectos Estratégicos del Sector Energético disponible en el portal \href{https://ventanillaunica.energia.gob.mx/}{\texttt{enlace}}, habilitada como un módulo para el seguimiento de proyectos del sector.\marcaSegunda{}\referenciaSegunda{}}


\subsubsection*{3. Reglas generales y de interpretación}


\paragraph*{3.1}
\mbox{}\par
\nopagebreak[4]


\puntoConvocatoria{I. Para efectos de la Convocatoria de Proyectos Estratégicos, los términos definidos en el Glosario aplican en todos los casos y pueden utilizarse en singular o plural, según el contexto.\referenciaBase{}}

\puntoConvocatoria{II. Las fechas previstas deben entenderse dentro del año de publicación de la Convocatoria de Proyectos Estratégicos y los días deben ser considerados como hábiles en términos de lo previsto en el artículo 28 de la Ley Federal de Procedimiento Administrativo, salvo que se indique expresamente que son naturales.\referenciaBase{}}

\puntoConvocatoria{III. El otorgamiento de un permiso en términos de la presente Convocatoria de Proyectos Estratégicos se encuentra sujeto al cumplimiento de los requisitos establecidos en la LSE, el RLSE, las Disposiciones de Permisos, las Disposiciones de Integración de SAEE y demás normatividad aplicable, así como lo previsto en esta Convocatoria.\referenciaBase{}}

\puntoConvocatoria{IV. Para efectos administrativos, corresponde a la SENER interpretar y resolver lo relativo a la Convocatoria de Proyectos Estratégicos, con base en criterios de certeza, competencia técnica y de planeación vinculante, con apoyo de la CNE para lo concerniente al otorgamiento de permisos de generación y del CENACE, en lo que respecta a los estudios de interconexión y conexión, contratos de interconexión y conexión, presentación y aplicación de garantías financieras, entre otros aspectos técnicos, así como de la CFE exclusivamente en materia de Proyectos en Asociación con la CFE cuando la persona solicitante opte por dicha modalidad.\referenciaBase{}}

\puntoConvocatoria{V. Los documentos que se presenten deben estar en idioma español; en caso de documentos en un idioma distinto, se debe anexar la traducción simple al español.\referenciaBase{}}

\puntoConvocatoria{VI. La capacidad de generación de energía eléctrica debe expresarse en MW con hasta dos decimales. La Potencia SAEE debe expresarse en MW y la Capacidad de Almacenamiento. Cuando exista diferencia entre número y letra, debe prevalecer la cifra numérica.\referenciaBase{}}

\puntoConvocatoria{VII. La fecha de Entrada en Operación Comercial de los proyectos debe ser conforme a lo establecido en el numeral 18; sin embargo, los proyectos pueden iniciar su operación comercial de manera anticipada, previa opinión favorable del CENACE, siempre que cumplan con los requisitos establecidos en la normatividad aplicable.\referenciaBase{}}

\puntoConvocatoria{VIII. Quedan a salvo los derechos de aquellas personas solicitantes cuya solicitud se deseche o no obtengan un permiso conforme a la Convocatoria de Proyectos Estratégicos, para iniciar los trámites correspondientes conforme a lo establecido en la LSE, el RLSE y demás normatividad aplicable. En cualquier caso, las personas interesadas deben cumplir con los requisitos aplicables y la planeación vinculante.\referenciaBase{}}

\puntoConvocatoria{IX. Los registros de manifestación de interés que se presenten después del plazo previsto en la Convocatoria de Proyectos Estratégicos, sin excepción alguna, no pueden continuar participando en las etapas subsecuentes de la misma.\referenciaBase{}}

\puntoConvocatoria{X. En el caso de que la persona solicitante haya ingresado una solicitud para la obtención de un permiso ante la CNE, previo a la emisión de la Convocatoria de Proyectos Estratégicos, puede solicitar que esta sea sujeta al procedimiento previsto en la Convocatoria. En este supuesto, la persona solicitante debe presentar el desistimiento de la o las solicitudes previas, en el caso de que se le otorgue el permiso correspondiente, en términos de esta Convocatoria.\referenciaBase{}}

\puntoConvocatoria{XI. El CENACE no debe considerar proyectos que correspondan a regiones donde no exista Capacidad Requerida para la Convocatoria de Proyectos Estratégicos o que esta haya sido cubierta.\marcaSegunda{}\referenciaSegunda{}}

\puntoConvocatoria{XII. Los proyectos asociados a solicitudes ingresadas al amparo de la presente Convocatoria pueden participar simultáneamente en otras convocatorias.\marcaSegunda{}\referenciaSegunda{}}

\puntoConvocatoria{XIII. Los estudios de conexión o interconexión, así como las obras de conexión o interconexión para cada proyecto aplican, según corresponda, según su condición operativa en el SEN.\referenciaBase{}}

\puntoConvocatoria{XIV. Para todo lo no previsto en la Convocatoria de Proyectos Estratégicos, se aplica lo dispuesto en la Ley Federal de Procedimiento Administrativo y demás normatividad del sector eléctrico vigente y aplicable.\referenciaBase{}}


\subsubsection*{4. Ventanilla Única de Proyectos Estratégicos del Sector Energético}


\paragraph*{4.1}
\mbox{}\par
\nopagebreak[4]


La VUPE constituye el medio por el cual se presentan los registros de manifestación de interés, se integran los expedientes de los proyectos que sigan el procedimiento previsto en la Convocatoria de Proyectos Estratégicos, siendo esta el medio principal para el desarrollo y seguimiento del proceso relativo a la propia Convocatoria.\referenciaBase{}


\paragraph*{4.2}
\mbox{}\par
\nopagebreak[4]


Las personas solicitantes que participen en las distintas etapas de la Convocatoria de Proyectos Estratégicos aceptan que a través de la VUPE se pueden presentar promociones, así como oír y recibir notificaciones relacionadas con la misma. De igual forma, aceptan que la OPE es el medio a través del cual deben presentar su solicitud de permiso y demás trámites asociados al permiso que, en su caso, les sea otorgado por parte de la CNE.\referenciaBase{}


\paragraph*{4.3}
\mbox{}\par
\nopagebreak[4]


A través de la VUPE, las personas solicitantes deben realizar las siguientes actividades:\referenciaBase{}


\puntoConvocatoria{I. Presentar su registro de manifestación de interés en participar en la misma y, en su caso, realizar uno o más proyectos a que se refiere el numeral 1.2, fracciones I y II de esta Convocatoria.\referenciaBase{}}

\puntoConvocatoria{II. Presentar la solicitud de estudios de conexión e interconexión y la presentación de garantías financiera y, en su caso, tipo de Proyecto en Asociación con la CFE.\referenciaBase{}}

\puntoConvocatoria{III. Consultar la información general sobre el avance de los trámites vinculados al proyecto respecto del cual solicitan un permiso, a partir de la progresiva interoperabilidad de la VUPE con las plataformas electrónicas de las distintas autoridades del sector que participan en las distintas etapas de la Convocatoria de Proyectos Estratégicos, lo que permite fortalecer la coordinación, transparencia y trazabilidad en el seguimiento prioritario de los proyectos registrados. Lo anterior no impide que dichas autoridades continúen utilizando los sistemas electrónicos y plataformas digitales con los que cuentan para el ejercicio de sus atribuciones.\referenciaBase{}}

\puntoConvocatoria{IV. En su caso, informar sobre el avance administrativo del proyecto, a partir del reporte del estado de los trámites asociados, que han presentado ante diversas autoridades federales, estatales y municipales.\referenciaBase{}}


\subsubsection*{5. Participación de las distintas autoridades en la Convocatoria de Proyectos Estratégicos y consideraciones sobre sus trámites}


\paragraph*{5.1}
\mbox{}\par
\nopagebreak[4]


La SENER, la CNE, el CENACE y la CFE, en el ámbito de sus atribuciones, deben realizar las gestiones necesarias para atender las solicitudes, estudios y procedimientos conforme a los plazos de la presente Convocatoria de Proyectos Estratégicos.\marcaSegunda{}\referenciaSegunda{}


\paragraph*{5.2}
\mbox{}\par
\nopagebreak[4]


Los trámites relacionados con los estudios de conexión e interconexión que se hayan iniciado previamente o se reciban durante las diversas etapas de la Convocatoria de Proyectos Estratégicos continúan sujetos a los plazos y términos previstos en la normatividad aplicable y en los instrumentos de planeación vinculante correspondientes, sin que la Convocatoria de Proyectos Estratégicos suspenda su tramitación, ni la atención de los demás proyectos.\referenciaBase{}


\paragraph*{5.3}
\mbox{}\par
\nopagebreak[4]


Los proyectos previamente ingresados o en trámite ante el CENACE o la CNE pueden ser reconocidos para efectos de la Convocatoria de Proyectos Estratégicos cuando las personas interesadas ratifiquen su solicitud o proyecto conforme a los mecanismos previstos en esta Convocatoria y se encuentren alineados con los instrumentos de planeación vinculante aplicables a la generación y al almacenamiento de energía eléctrica. En esos casos, el CENACE puede priorizar la realización, actualización o revalidación de los estudios correspondientes.\referenciaBase{}


Si la información presentada no es suficiente para revalidar o actualizar estudios de conexión e interconexión ya existentes, el CENACE debe notificarlo a la persona interesada por medio de la VUPE y señalar, en su caso, los ajustes que resulten procedentes.\referenciaBase{}


\paragraph*{5.4}
\mbox{}\par
\nopagebreak[4]


Los estudios realizados por el CENACE como parte del procedimiento previsto en esta Convocatoria únicamente son válidos para la misma, por lo que, en caso de que alguna persona solicitante pretenda utilizar, en la vía ordinaria, los estudios de conexión e interconexión realizados al amparo de dicho procedimiento respecto de proyectos a los que no se les otorgue un permiso, deben presentarse al CENACE, a efecto de que defina la vigencia o no de los mismos.\referenciaBase{}


A través de la VUPE, el CENACE debe notificar a las personas solicitantes, respecto de los proyectos a los que deba realizar los estudios correspondientes o bien, de aquellos cuyos estudios hayan sido revalidados.\referenciaBase{}


\subsubsection*{6. Etapas y Actividades}


\paragraph*{6.1}
\mbox{}\par
\nopagebreak[4]


Las personas interesadas en desarrollar uno o más proyectos a los que se refiere el numeral 1.2, fracciones I y II, deben llevar a cabo las actividades descritas, en las etapas y plazos previstos en el numeral 17 de la Convocatoria para Proyectos Estratégicos.\referenciaBase{}


\paragraph*{6.2}
\mbox{}\par
\nopagebreak[4]


Las etapas y actividades previstas en este numeral aplican según el proyecto se desarrolle o no, como Proyecto en Asociación con la CFE.\referenciaBase{}


Las etapas y plazos previstos en esta Convocatoria son secuenciales y su desarrollo se encuentra sujeto al cumplimiento de cada etapa inmediata anterior correspondiente, por lo que en caso de no cumplirse alguna de ellas, la persona solicitante no puede continuar con las etapas subsecuentes del procedimiento, salvo aquellas etapas que determinen una excepción.\marcaTercera{}\referenciaSegunda{}


Lo anterior, sin perjuicio de los supuestos específicos de conclusión o desechamiento expresamente previstos en esta Convocatoria.\marcaTercera{}\referenciaSegunda{}


\subsubsection*{7. Estudios de conexión e interconexión y pagos correspondientes}


\paragraph*{7.1}
\mbox{}\par
\nopagebreak[4]


Para efectos de la Convocatoria de Proyectos Estratégicos, los estudios de interconexión y conexión a realizar, según corresponda, son el estudio de impacto y el estudio de instalaciones, sin que resulte necesaria la realización del estudio indicativo.\referenciaBase{}


\paragraph*{7.2}
\mbox{}\par
\nopagebreak[4]


Los pagos para los estudios de interconexión de los proyectos de generación de energía eléctrica que correspondan, están sujetos a lo previsto en el MIC, como se estipula en su Anexo I, Tabla de Costos de los Estudios del CENACE, numeral I.1 Estudios de Interconexión.\referenciaBase{}


\paragraph*{7.3}
\mbox{}\par
\nopagebreak[4]


Para el caso de los pagos para los estudios de conexión e interconexión de proyectos de Sistemas de Almacenamiento de Energía Eléctrica, el trato que se les debe dar a dichos sistemas es de Centros de Carga cuando retiren energía del SEN y de Centrales Eléctricas cuando inyecten energía al SEN, en lo que resulte aplicable en términos del MIC.\referenciaBase{}


\paragraph*{7.4}
\mbox{}\par
\nopagebreak[4]


Cuando exista un pago previo relacionado con estudios de conexión o interconexión que la persona interesada pretenda acreditar para efectos de la Convocatoria de Proyectos Estratégicos, debe precisarlo en la VUPE para que el CENACE determine si procede su reconocimiento o actualización.\referenciaBase{}


En caso de que no exista un pago previo o la información resulte insuficiente, la persona interesada debe cubrir el pago que determine el CENACE para los estudios de conexión e interconexión que correspondan.\referenciaBase{}


\paragraph*{7.5}
\mbox{}\par
\nopagebreak[4]


Cuando exista una modificación técnica entre la información presentada en la VUPE y la previamente ingresada ante el CENACE, este debe determinar los estudios que deban realizarse nuevamente.\referenciaBase{}


Cuando el proyecto cuente con estudios de conexión e interconexión cuya vigencia haya concluido, el CENACE debe determinar, conforme a la normatividad aplicable, si procede su actualización o la realización total o parcial de nuevos estudios.\referenciaBase{}


\paragraph*{7.6}
\mbox{}\par
\nopagebreak[4]


La persona interesada debe atender los requerimientos de información técnica adicional que formule el CENACE para definir el alcance de los estudios de conexión e interconexión correspondientes.\referenciaBase{}


\paragraph*{7.7}
\mbox{}\par
\nopagebreak[4]


Para efectos de la Convocatoria de Proyectos Estratégicos, la procedencia, vigencia, actualización, revalidación y reconocimiento de estudios se determinan conforme a las Disposiciones de Integración de SAEE.\referenciaBase{}


\subsubsection*{8. Requisito de aceptación de Obras de Interconexión, Obras de Conexión y Obras de Refuerzo}


\paragraph*{8.1}
\mbox{}\par
\nopagebreak[4]


Una vez que el CENACE haya notificado los resultados de los estudios de conexión o interconexión, así como los costos de las Obras de Interconexión, Conexión u Obras de Refuerzo que correspondan, las personas solicitantes que decidan continuar en el procedimiento deben presentar el Formato de Aceptación de Obras dentro del plazo previsto en la Convocatoria de Proyectos Estratégicos.\referenciaBase{}


Presentado el formato al que hace referencia el párrafo anterior y celebrado el contrato o instrumento jurídico que formalice el Proyecto en Asociación con la CFE en términos de la normatividad aplicable, la persona solicitante o la CFE pueden presentar la solicitud de permiso correspondiente para que la CNE realice su evaluación de conformidad con los requisitos establecidos en la LSE, el RLSE, la LPTE, el RLPTE, la Convocatoria de Proyectos Estratégicos y demás disposiciones aplicables.\marcaSegunda{}\referenciaSegunda{}


\paragraph*{8.2}
\mbox{}\par
\nopagebreak[4]


En su caso, el CENACE debe determinar si procede la revalidación, actualización o la realización de nuevos estudios conforme a las Disposiciones de Integración de SAEE y demás normatividad aplicable.\referenciaBase{}


El CENACE, a través de la VUPE, debe hacer del conocimiento de cada persona solicitante por medio de un solo documento los resultados de los estudios de conexión e interconexión, indicándose el costo estimado de las Obras de Refuerzo, así como las Obras de Interconexión y Conexión que se deben cubrir en caso de que se le otorgue el permiso correspondiente en términos de la Convocatoria de Proyectos Estratégicos. En el caso de Sistemas de Almacenamiento de Energía Eléctrica, se estará a lo previsto en la normatividad aplicable.\referenciaBase{}


\paragraph*{8.3}
\mbox{}\par
\nopagebreak[4]


Los resultados de los estudios de conexión e interconexión que realice y notifique el CENACE deben indicar al menos lo siguiente:\referenciaBase{}


\puntoConvocatoria{I. Descripción detallada de las Obras de Interconexión o Conexión y Obras de Refuerzo necesarias en la Red Nacional de Transmisión o las Redes Generales de Distribución, así como el Nivel de Tensión en el Punto de Interconexión.\referenciaBase{}}


\puntoConvocatoria{II. Costo estimado, en pesos mexicanos, de las Obras de Refuerzo. Cuando estas sean comunes a dos o más proyectos, los costos se asignan de manera proporcional a la potencia en MW de cada proyecto de generación o almacenamiento de energía eléctrica, considerando que se cubra el requerimiento de Capacidad Requerida para la Convocatoria de Proyectos Estratégicos, siempre que la o las personas solicitantes acepten los términos y continúen con lo previsto en la Convocatoria de Proyectos Estratégicos.\marcaSegunda{}\referenciaSegunda{}}


\puntoConvocatoria{III. Observaciones técnicas, restricciones operativas y condiciones específicas de conexión e interconexión.\referenciaBase{}}


Sobre los estudios de conexión e interconexión realizados por el CENACE al amparo de la Convocatoria de Proyectos Estratégicos no proceden aclaraciones, sin perjuicio de lo que prevea la normatividad aplicable respecto de su actualización o reevaluación.\referenciaBase{}


Si la persona solicitante decide continuar en el proceso, debe presentar a través de la VUPE, a más tardar en el plazo previsto en el numeral 17 de la Convocatoria de Proyectos Estratégicos, el Formato Aceptación de Obras, en el que acepte la ejecución de las Obras de Refuerzo y Obras de Interconexión necesarias. Lo anterior es una condición necesaria para continuar el proceso y obtener el permiso solicitado.\referenciaBase{}


En caso de no recibirse esta aceptación a través de la VUPE, en los términos previstos, el CENACE debe emitir el oficio correspondiente en el que se notifique lo conducente.\marcaSegunda{}\referenciaSegunda{}


En el supuesto de que los proyectos de dos o más personas solicitantes hayan coincidido en relación con la Gerencia de Control Regional y que, con base en ello, se hayan determinado las Obras de Refuerzo notificadas a dichas personas y que alguna de estas no las haya aceptado, el CENACE debe reevaluar por única ocasión los resultados de los estudios, así como el costo de Obras de Refuerzo y las Obras de Interconexión concernientes a las personas solicitantes que sí habían aceptado el costo de las Obras de Refuerzo originales.\referenciaBase{}


El CENACE debe notificar la información derivada de tal reevaluación dentro de los 5 días hábiles siguientes a que se actualice dicho supuesto, a efecto de que las personas solicitantes correspondientes puedan aceptar las nuevas características de las referidas obras y los nuevos costos. Las personas solicitantes a quienes se notifique la reevaluación deben presentar, a través de la VUPE, la aceptación expresa de las obras y costos redefinidos dentro de los 5 días hábiles siguientes a la notificación del CENACE, conforme a lo establecido en el calendario previsto en la presente Convocatoria.\referenciaBase{}


Una vez aceptadas las características y los costos asociados a las Obras de Interconexión y Obras de Refuerzo, es irrevocable su ejecución, por lo que cualquier negativa o incumplimiento posterior implica el incumplimiento de los términos u obligaciones contenidos en el permiso otorgado al amparo de la Convocatoria de Proyectos Estratégicos, del contrato o instrumento correspondiente que se celebre, en su caso, así como la liberación de la capacidad asignada.\referenciaBase{}


\paragraph*{8.4}
\mbox{}\par
\nopagebreak[4]


La persona solicitante debe considerar que, además de los costos asociados a las Obras de Refuerzo que resulten de los estudios de conexión e interconexión realizados por el CENACE, debe cubrir el costo relacionado con la infraestructura que integra su proyecto con la Red Nacional de Transmisión o las Redes Generales de Distribución en el Punto de Interconexión, de conformidad con la normatividad aplicable.\referenciaBase{}


\subsubsection*{9. Alineación a la planeación vinculante}


\paragraph*{9.1}
\mbox{}\par
\nopagebreak[4]


La CNE debe corroborar que se cuente con los elementos necesarios para determinar que los proyectos participantes cumplan con los criterios de planeación vinculante que se requieren para el otorgamiento de un permiso bajo la Convocatoria de Proyectos Estratégicos; para ello, debe solicitar la siguiente información:\referenciaBase{}


\puntoConvocatoria{I. A la SENER, la información relacionada con el monto de inversión destinado al Plan de Gestión Social, así como lo referente a la prevalencia, innovación y la contribución a la mitigación de gases de efecto invernadero.\referenciaBase{}}


\puntoConvocatoria{II. Al CENACE, los estudios técnicos referidos en la Convocatoria de Proyectos Estratégicos, incluyendo la información necesaria para verificar el cumplimiento de los criterios de planeación vinculante, según le corresponda, establecidos en las Disposiciones de la Planeación Vinculante y dar continuidad a las etapas subsecuentes del procedimiento.\marcaSegunda{}\referenciaSegunda{}}


El CENACE debe proporcionar a la CNE la relación de los proyectos que proporcionen beneficios al SEN conforme a la evaluación que realice aplicando los criterios de las Disposiciones de la Planeación Vinculante, ordenados conforme al grado de beneficio que representen para el SEN.\referenciaBase{}


La SENER y el CENACE deben entregar a la CNE la información solicitada conforme a las etapas y plazos previstos en el numeral 17 de la Convocatoria de Proyectos Estratégicos.\referenciaBase{}


Los proyectos que no cumplan con la alineación a la planeación vinculante prevista en las Disposiciones de la Planeación Vinculante correspondiente, no son elegibles para solicitar y, en su caso, obtener un permiso al amparo de la Convocatoria de Proyectos Estratégicos.\referenciaBase{}


\subsubsection*{10. Evaluación de los proyectos por parte del Grupo de Análisis Técnico}


\paragraph*{10.1}
\mbox{}\par
\nopagebreak[4]


La SENER debe someter a consideración del Grupo de Análisis Técnico la evaluación de los proyectos que cumplan con los requisitos establecidos en la LSE, el RLSE, la LPTE, el RLPTE, la Convocatoria de Proyectos Estratégicos y demás disposiciones aplicables, para lo cual debe remitir los expedientes administrativos correspondientes.\marcaSegunda{}\referenciaSegunda{}


El Grupo de Análisis Técnico debe:\referenciaBase{}


\puntoConvocatoria{I. Emitir la opinión técnica a que se refiere la etapa 10 del numeral 17.1, respecto de la evaluación realizada por la SENER y el CENACE de los proyectos presentados al amparo de la Convocatoria de Proyectos Estratégicos.\marcaSegunda{}\referenciaSegunda{}}


\puntoConvocatoria{II. Identificar aquellos proyectos que proporcionen los mayores beneficios para el SEN y el Estado, en el caso de que los proyectos que resulten elegibles para obtener un permiso al amparo de la Convocatoria de Proyectos Estratégicos. Para el cumplimiento de lo anterior, deben considerar al menos lo siguiente respecto de los proyectos evaluados: Entrada en Operación Comercial, ubicación, elementos técnicos básicos del proyecto, alivio de congestión, aportación de potencia, mejora de la confiabilidad regional y del SEN, flexibilidad operativa, aportaciones socioambientales, tecnología y desempeño proyectado del Sistema de Almacenamiento de Energía Eléctrica, forma de ejecución y financiamiento, y eficiencia económica.\marcaSegunda{}\referenciaSegunda{}}


\paragraph*{10.2}
\mbox{}\par
\nopagebreak[4]


La SENER debe notificar la opinión del GAT a la CFE, a la CNE y a las personas solicitantes a través de la VUPE, para que, en su caso, se continúe con lo previsto en el numeral 17.3 de esta Convocatoria.\marcaSegunda{}\referenciaSegunda{}


\paragraph*{10.3}
\mbox{}\par
\nopagebreak[4]


Una vez emitida la opinión del GAT, en caso de no continuar con el procedimiento, las personas solicitantes pueden presentar por las vías ordinarias las solicitudes para la obtención de los diversos permisos, autorizaciones y demás actos administrativos correspondientes, en los términos previstos en la normatividad aplicable.\marcaSegunda{}\referenciaSegunda{}


\subsubsection*{11. Procedencia de Proyectos en Asociación con la CFE}


\paragraph*{11.1}
\mbox{}\par
\nopagebreak[4]


Una vez emitida la opinión técnica por parte del Grupo de Análisis Técnico, la SENER debe proporcionar a la CFE el expediente del proyecto, a través de la VUPE, para que, en su caso, se continúe el procedimiento conforme a lo previsto en el numeral 17.3 de la presente Convocatoria.\marcaSegunda{}\referenciaSegunda{}


La remisión del expediente a la CFE no implica la aprobación del proyecto, del permiso correspondiente, ni del esquema de asociación, y queda sujeta a la resolución que emitan las autoridades competentes y la CFE conforme a la normatividad aplicable.\referenciaBase{}


\subsubsection*{12. Aprobación de permisos y emisión de los títulos de permisos}


\paragraph*{12.1}
\mbox{}\par
\nopagebreak[4]


La CNE debe resolver sobre la solicitud de permiso presentada por la persona solicitante, considerando los plazos establecidos en el numeral 17.4 de la presente Convocatoria y de conformidad con lo previsto en la LSE, el RLSE y la demás normatividad aplicable, para lo cual le debe ser remitido el expediente administrativo que incluya toda la información presentada por la persona solicitante, así como la proporcionada por la SENER, el CENACE y la CFE.\marcaSegunda{}\referenciaSegunda{}


\paragraph*{12.2}
\mbox{}\par
\nopagebreak[4]


En los casos en que se presente la solicitud de permiso correspondiente, la CNE debe verificar la siguiente información:\referenciaBase{}


\begin{longtable}{@{}>{\raggedright\arraybackslash}p{0.23\textwidth}>{\raggedright\arraybackslash}p{0.34\textwidth}>{\raggedright\arraybackslash}p{0.39\textwidth}@{}}
\toprule
\senerth{No.} & \senerth{Concepto} & \senerth{Cumple} \\
\midrule
\endfirsthead
\toprule
\senerth{No.} & \senerth{Concepto} & \senerth{Cumple} \\
\midrule
\endhead
1 & La persona solicitante manifestó su interés en obtener un permiso en términos de la Convocatoria de Proyectos Estratégicos.\referenciaBase{} & Sí / No \\
2 & La persona solicitante presentó la evidencia del pago de los estudios del CENACE referidos en la Convocatoria de Proyectos Estratégicos, esto es, la evidencia de los pagos anteriormente realizados, o bien, la de pagos nuevos.\referenciaBase{} & Sí / No \\
3 & La solicitud presentada corresponde a un proyecto que se alinea con lo establecido en el numeral 17 de la Convocatoria de Proyectos Estratégicos.\referenciaBase{} & Sí / No \\
4 & La persona solicitante cumple con los requisitos establecidos en las Disposiciones de Permisos.\referenciaBase{} & Sí / No \\
5 & La persona solicitante presentó la aceptación de las Obras de Refuerzo y de interconexión correspondientes conforme al Formato Aceptación de Obras.\referenciaBase{} & Sí / No \\
6 & En el caso de Proyectos con Asociación de la CFE, contar con el instrumento que acredite la asociación con la CFE.\referenciaBase{} & Sí / No \\
\bottomrule
\end{longtable}


El incumplimiento de cualquiera de los requisitos antes señalados tiene como consecuencia el no otorgamiento de un permiso con base en la Convocatoria de Proyectos Estratégicos.\referenciaBase{}


Las solicitudes que se sometan a la aprobación del órgano competente de la CNE deben cumplir con los requisitos establecidos en la LSE, el RLSE, la LPTE, el RLPTE y demás disposiciones aplicables, incluyendo lo previsto en la Convocatoria de Proyectos Estratégicos. Asimismo, los proyectos presentados se deben alinear a la planeación vinculante y, en su caso, ofrecer las mejores condiciones para el SEN.\referenciaBase{}


\paragraph*{12.3}
\mbox{}\par
\nopagebreak[4]


Los títulos de permiso se emiten, según corresponda, en términos de la LSE, el RLSE y las Disposiciones de Permisos.\referenciaBase{}


\paragraph*{12.4}
\mbox{}\par
\nopagebreak[4]


La CNE debe emitir y notificar los títulos de permisos que cumplan con lo previsto en la Convocatoria de Proyectos Estratégicos y sean aprobados por su órgano competente. Dichos permisos deben contener, al menos, los siguientes plazos y términos que las personas permisionarias deben cumplir:\referenciaBase{}


\puntoConvocatoria{I. Presentar la documentación que acredite el esquema de financiamiento del proyecto en un plazo máximo de 8 meses contados a partir de la notificación del permiso o instrumento correspondiente.\referenciaBase{}}

\puntoConvocatoria{II. Acreditar, previo a la firma del contrato de conexión o interconexión correspondiente, la presentación de la garantía financiera por las Obras de Interconexión o Conexión u Obras de Refuerzo resultado del estudio que hayan aceptado ante el CENACE.\referenciaBase{}}

\puntoConvocatoria{III. Obtener la autorización definitiva en materia de Manifestación de Impacto Social del Sector Energético o, en su caso, Evaluación del Impacto Social, así como la autorización en materia de impacto ambiental emitida por la autoridad competente, y cualquier documento que lo modifique o sustituya para iniciar las obras correspondientes, en un plazo máximo de 6 meses contados a partir de la notificación del permiso o instrumento correspondiente, salvo en los casos en los que sea necesaria la realización de una consulta previa, libre e informada.\referenciaBase{}}


\subsubsection*{13. Presentación y devolución de las garantías financieras asociadas a contratos de conexión e interconexión derivados de los permisos otorgados bajo la Convocatoria de Proyectos Estratégicos}


El contenido íntegro del numeral 13 forma parte de la estructura original de la Convocatoria y debe conservarse con la redacción publicada el 15 de mayo de 2026, al no haberse identificado modificación expresa en los acuerdos posteriores adjuntos.\referenciaBase{}\referenciaSegunda{}


\subsubsection*{14. Publicación periódica de capacidades y resultados de la Convocatoria de Proyectos Estratégicos}


La SENER puede publicar y, en su caso, actualizar la Capacidad Requerida para la Convocatoria de Proyectos Estratégicos, así como el listado de proyectos que cumplieron con los requisitos correspondientes y con los criterios de planeación vinculante, y que obtuvieron un permiso con base en esta Convocatoria, en la página \href{https://ventanilla.energia.gob.mx/}{\texttt{enlace}}.\marcaSegunda{}\referenciaSegunda{}


Las publicaciones a las que hace referencia este numeral tienen carácter informativo y no sustituyen los actos, estudios, resoluciones o notificaciones que correspondan conforme a la normatividad aplicable.\marcaSegunda{}\referenciaSegunda{}


\subsubsection*{15. Suscripción del contrato de conexión o interconexión correspondiente a los permisos otorgados bajo la Convocatoria de Proyectos Estratégicos}


El numeral 15 forma parte del texto original del instrumento y no aparece entre los numerales expresamente modificados por los acuerdos posteriores adjuntos, por lo que debe conservarse con la redacción de la Convocatoria original para una reproducción íntegra del documento.\referenciaBase{}\referenciaSegunda{}


\subsubsection*{16. Seguimiento de permisos hasta la Entrada en Operación Comercial de los proyectos}


El numeral 16 forma parte del texto original del instrumento y no aparece entre los numerales expresamente modificados por los acuerdos posteriores adjuntos, por lo que debe conservarse con la redacción de la Convocatoria original para una reproducción íntegra del documento.\referenciaBase{}\referenciaSegunda{}


\subsubsection*{17. Etapas, plazos y actividades de la Convocatoria de Proyectos Estratégicos\marcaSegunda{}}


\paragraph*{17.1}
\mbox{}\par
\nopagebreak[4]


Para el desarrollo de la Convocatoria de Proyectos Estratégicos, el procedimiento de registro de manifestación de interés y solicitud de estudios de conexión e interconexión queda sujeto a las etapas, plazos y actividades siguientes.\marcaSegunda{}\referenciaSegunda{}


En todos los casos, el procedimiento queda sujeto a lo siguiente:\marcaSegunda{}\referenciaSegunda{}


\begin{longtable}{@{}>{\raggedright\arraybackslash}p{0.17\textwidth}>{\raggedright\arraybackslash}p{0.24\textwidth}>{\raggedright\arraybackslash}p{0.25\textwidth}>{\raggedright\arraybackslash}p{0.29\textwidth}@{}}
\toprule
\senerth{No. de etapa} & \senerth{Alcance} & \senerth{Fechas/plazos} & \senerth{Descripción de la actividad} \\
\midrule
\endfirsthead
\toprule
\senerth{No. de etapa} & \senerth{Alcance} & \senerth{Fechas/plazos} & \senerth{Descripción de la actividad} \\
\midrule
\endhead
1 & Publicación de la Convocatoria de Proyectos Estratégicos & 15 de mayo & SENER publica en el DOF la Convocatoria de Proyectos Estratégicos.\referenciaSegunda{} \\
2 & Registro de manifestación de interés y solicitud de estudios de interconexión y conexión & 2 de junio al 2 de septiembre de 2026 & La persona solicitante debe registrar su manifestación de interés y la solicitud de estudios de conexión e interconexión, según corresponda, a través de la VUPE en el formato previsto para tales efectos.\marcaPrimera{} \marcaSegunda{}\referenciaPrimera{}\referenciaSegunda{} \\
3 & Prevención y atención de prevención de registro y solicitud & 7 días & Una vez recibida la solicitud de estudios de conexión e interconexión, según corresponda, cuando la solicitud no contenga los datos e información correspondientes o no cumpla con los requisitos establecidos, el CENACE debe desahogar la etapa de prevención dentro de un plazo de siete días hábiles. Para tal efecto, debe prevenir a la persona solicitante dentro de los primeros dos días hábiles de esta etapa, para que subsane lo conducente dentro de los 5 días hábiles siguientes.\referenciaSegunda{} \\
 &  &  & Vencido el plazo previsto en el párrafo anterior, en caso de que no se atienda la prevención en tiempo y forma, el CENACE debe desechar la solicitud correspondiente y notificar dicha determinación a la persona solicitante.\referenciaSegunda{} \\
4 & CFE manifiesta si es de su interés que se evalúe el proyecto presentado para asociarse en términos de sus directrices y de la Convocatoria de Proyectos Estratégicos & 10 días & La CFE debe revisar los registros de manifestación de interés de los proyectos presentados en la Convocatoria de Proyectos Estratégicos, para manifestar si es de su interés que se evalúe el proyecto presentado para asociarse en términos de sus directrices y de la Convocatoria de Proyectos Estratégicos, e informarlo a través de la VUPE, para continuar con el procedimiento previsto en esta Convocatoria.\referenciaSegunda{} \\
 &  &  & En el caso de que la CFE manifieste que no es de su interés que se evalúe el proyecto, se da por concluido el procedimiento previsto en esta Convocatoria.\referenciaSegunda{} \\
5 & Determinación de pagos por estudios de conexión e interconexión & 2 días & Determinados los pagos por los estudios correspondientes, por medio de la VUPE, el CENACE debe comunicar a las personas solicitantes alguna de las siguientes opciones: 1. El monto a pagar, para lo cual se le notifica a la persona solicitante la línea de captura de pago correspondiente, o 2. Que la información presentada resulta suficiente para la revalidación de los estudios correspondientes, por lo que no debe realizar un pago adicional. En este supuesto, no aplican las etapas 6 y 7 subsecuentes.\referenciaSegunda{} \\
6 & Pago de estudios de conexión e interconexión & 5 días & Cuando así proceda, la persona solicitante debe presentar, por medio de la VUPE, el comprobante del pago de los estudios de conexión e interconexión que, en su caso, haya informado el CENACE.\referenciaSegunda{} \\
7 & Realización de estudios de interconexión y conexión & 30 días & Una vez validado el pago por los estudios, el CENACE debe realizar los estudios correspondientes conforme a lo previsto en el MIC. La persona solicitante debe revisar las notificaciones que se realicen a través de la VUPE.\referenciaSegunda{} \\
8 & Notificación de resultados de estudios y costos de Obras de Refuerzo y costos de interconexión & 1 día & El CENACE debe notificar a la persona solicitante y a la CNE el resultado de los estudios de conexión e interconexión o su revalidación, según corresponda, a más tardar el día hábil siguiente transcurrido el plazo para la elaboración de los estudios.\referenciaSegunda{} \\
9 & Aceptación de costos de Obras de Refuerzo y de Interconexión & 5 días & Una vez notificado el resultado de los estudios o su revalidación, la persona solicitante, en su caso, debe presentar a través de la VUPE, el Anexo 19.3, Formato Aceptación de Obras. En el caso de que la persona solicitante no acepte los costos de las Obras de Refuerzo y de Interconexión, el CENACE debe concluir la solicitud de interconexión y conexión y, por lo tanto, el procedimiento previsto en esta Convocatoria.\referenciaSegunda{} \\
10 & Revisión de información técnica y emisión de Opinión del Grupo de Análisis Técnico & 5 días & El Grupo de Análisis Técnico debe revisar la información técnica que, en su caso, le remita la SENER, el CENACE o la CFE, respecto de los proyectos que presenten su registro de manifestación de interés en términos de la Convocatoria de Proyectos Estratégicos y emitir su opinión técnica.\referenciaSegunda{} \\
\bottomrule
\end{longtable}


\paragraph*{17.2}
\mbox{}\par
\nopagebreak[4]


La SENER debe notificar, a través de la VUPE, la opinión emitida por el Grupo de Análisis Técnico a la que se refiere la etapa 10 prevista en el numeral 17.1 de la Convocatoria de Proyectos Estratégicos, a la persona solicitante, a la CFE y a la CNE.\marcaSegunda{}\referenciaSegunda{}


\paragraph*{17.3}
\mbox{}\par
\nopagebreak[4]


Emitida la opinión del Grupo de Análisis Técnico a la que se refiere la etapa 10 prevista en el numeral 17.1, la CFE debe iniciar el procedimiento correspondiente para que, en su caso, formalice el Proyecto en Asociación con la CFE.\marcaSegunda{}\referenciaSegunda{}


Para la formalización de los contratos o instrumentos jurídicos del Proyecto Estratégico con la CFE, se debe estar a lo previsto en la LSE, el RLSE, la LPTE, el RLPTE, la Ley de la Empresa Pública del Estado, Comisión Federal de Electricidad y su reglamento y, en su caso, a los Lineamientos de los Esquemas para el Desarrollo Mixto de la Empresa Pública del Estado, Comisión Federal de Electricidad y demás normatividad aplicable.\marcaTercera{}\referenciaSegunda{}


En caso de que no se formalice el Proyecto en Asociación con la CFE, la empresa pública del Estado lo debe informar a la SENER a efecto de que esta notifique, a través de la VUPE, la conclusión del procedimiento. Lo anterior, sin perjuicio de que la persona interesada pueda participar en otras convocatorias, conforme a la normatividad aplicable.\marcaTercera{}\referenciaSegunda{}


\paragraph*{17.4}
\mbox{}\par
\nopagebreak[4]


Formalizado el Proyecto en Asociación con la CFE, la persona solicitante o la CFE debe presentar la solicitud del permiso que corresponda de conformidad con las etapas, plazos y actividades siguientes:\marcaSegunda{}\referenciaSegunda{}


\begin{longtable}{@{}>{\raggedright\arraybackslash}p{0.17\textwidth}>{\raggedright\arraybackslash}p{0.24\textwidth}>{\raggedright\arraybackslash}p{0.25\textwidth}>{\raggedright\arraybackslash}p{0.29\textwidth}@{}}
\toprule
\senerth{No. de etapa} & \senerth{Alcance} & \senerth{Plazos} & \senerth{Descripción de actividad} \\
\midrule
\endfirsthead
\toprule
\senerth{No. de etapa} & \senerth{Alcance} & \senerth{Plazos} & \senerth{Descripción de actividad} \\
\midrule
\endhead
1 & Presentación de la solicitud de permiso & 15 días & La persona solicitante o la CFE debe ingresar, a través de la OPE, la solicitud de permiso correspondiente, de conformidad con lo establecido en las Disposiciones de Permisos. En su caso, la persona solicitante puede optar por alguna de las siguientes opciones, según corresponda: 1. Sin trámite previo. Cuando la persona solicitante no cuente con ninguna solicitud de permiso en trámite ante la CNE, debe presentar el formato de solicitud de permiso correspondiente a través de la OPE, o 2. Con trámite previo. Cuando la persona solicitante ya cuente con una solicitud de permiso previa a la presentación del registro de manifestación de interés en términos de la Convocatoria de Proyectos Estratégicos, debe presentar a través de la OPE el Formato Solicitudes de Permiso con Trámite Previo. Durante esta etapa la SENER y el CENACE deben remitir a la CNE la información a la que se refiere el numeral 9.1 de esta Convocatoria.\referenciaSegunda{} \\
2 & Revisión de requisitos y prevención & 10 días & Recibida una solicitud de permiso, la CNE debe revisar que la misma cumple con lo previsto en las Disposiciones de Permisos. En su caso, la CNE debe realizar la prevención correspondiente cuando la solicitud no cuente con la información o requisitos correspondientes, a través de la OPE. La persona solicitante o la CFE debe dar atención en tiempo y forma a la prevención que, en su caso, se le haya realizado, a través de la OPE; de no hacerlo, se debe desechar la solicitud.\referenciaSegunda{} \\
3 & Evaluación de criterios de planeación vinculante por parte de la CNE & 7 días & La CNE debe llevar a cabo la evaluación conforme a lo previsto en las Disposiciones de la Planeación Vinculante.\referenciaSegunda{} \\
4 & Resolución de las solicitudes de permisos por parte de la CNE & 5 días & La CNE debe resolver lo conducente respecto de las solicitudes de permisos correspondientes.\referenciaSegunda{} \\
5 & Notificación de la resolución de las solicitudes de permisos & 1 día & La CNE debe notificar la resolución de la solicitud del permiso de que se trate, a través de la OPE.\referenciaSegunda{} \\
6 & Desistimiento de solicitudes de permisos & 4 días & Tratándose de solicitudes con trámite previo, la persona solicitante a la que se le haya notificado el otorgamiento del permiso correspondiente debe desistirse de la solicitud de permiso previa ante la CNE.\referenciaSegunda{} \\
7 & Ejecución y seguimiento del programa de obras & A partir de la notificación del título de permiso o instrumento correspondiente & La persona solicitante o la CFE debe cumplir con los plazos y términos establecidos en el título de permiso otorgado.\referenciaSegunda{} \\
8 & Presentación y validación de garantías financieras & Previo a la firma del contrato de interconexión & La persona solicitante o la CFE debe presentar las garantías financieras correspondientes para poder llevar a cabo la celebración del contrato para la interconexión, conexión o ambas a la Red Nacional de Transmisión o a las Redes Generales de Distribución, según corresponda. Presentadas las garantías financieras, el CENACE debe validarlas y, en su caso, instruir la firma del contrato de interconexión o conexión correspondiente.\referenciaSegunda{} \\
9 & Suscripción del contrato de interconexión & 10 días & La persona solicitante o la CFE debe presentar la documentación e información que requiere el CENACE para suscribir el contrato o instrumento correspondiente. Una vez presentada la garantía financiera correspondiente y cumplidos los requisitos aplicables, las entidades responsables y la persona solicitante, a través de la VUPE, deben realizar las gestiones necesarias para la suscripción del contrato para la interconexión, conexión o ambas, en términos de la normatividad aplicable.\referenciaSegunda{} \\
10 & Interconexión o conexión física al SEN & 10 días & La CFE debe realizar las acciones necesarias para realizar la interconexión o conexión del proyecto de que se trate al SEN.\referenciaSegunda{} \\
11 & Entrada en Operación Comercial & Según título de permiso & La persona solicitante y la CFE deben realizar las acciones necesarias para lograr la Entrada en Operación Comercial del proyecto. La SENER, a través del Consejo de Planeación Energética, así como la CNE y el CENACE, deben dar seguimiento al desarrollo de los proyectos hasta su Entrada en Operación Comercial.\referenciaSegunda{} \\
\bottomrule
\end{longtable}


\paragraph*{17.4.1}
\mbox{}\par
\nopagebreak[4]


Para efectos del procedimiento de otorgamiento de permisos en términos de la Convocatoria de Proyectos Estratégicos, los plazos y etapas para su atención corresponden a los señalados en la misma.\marcaSegunda{}\referenciaSegunda{}


\subsubsection*{18. Referencias de Capacidad Requerida para la Convocatoria de Proyectos Estratégicos de Sistemas de Almacenamiento de Energía Eléctrica y Generación de Energía Eléctrica\marcaSegunda{}}


La referencia de Capacidad Requerida para la Convocatoria de Proyectos Estratégicos establece las cantidades para identificar las necesidades generales de generación y almacenamiento de energía eléctrica, dando prioridad a los proyectos cuya Entrada en Operación Comercial esté prevista antes de 2030, de acuerdo con lo siguiente.\marcaSegunda{}\referenciaSegunda{}


\paragraph*{18.1}
\mbox{}\par
\nopagebreak[4]


Para la generación de energía eléctrica, la referencia de Capacidad Requerida para la Convocatoria de Proyectos Estratégicos considera la planeación vinculante y las condiciones operativas del SEN.\marcaSegunda{}\referenciaSegunda{}


\paragraph*{18.2}
\mbox{}\par
\nopagebreak[4]


\begin{longtable}{@{}>{\raggedright\arraybackslash}p{0.23\textwidth}>{\raggedright\arraybackslash}p{0.34\textwidth}>{\raggedright\arraybackslash}p{0.39\textwidth}@{}}
\toprule
\senerth{Gerencia de Control Regional} & \senerth{Capacidad Requerida para la Convocatoria de Proyectos Estratégicos (MW)\marcaSegunda{}\referenciaSegunda{}} & \senerth{Horas de almacenamiento} \\
\midrule
\endfirsthead
\toprule
\senerth{Gerencia de Control Regional} & \senerth{Capacidad Requerida para la Convocatoria de Proyectos Estratégicos (MW)\marcaSegunda{}\referenciaSegunda{}} & \senerth{Horas de almacenamiento} \\
\midrule
\endhead
Baja California & 140 & 3 horas \referenciaSegunda{} \\
Baja California Sur & 50 & 3 horas \referenciaSegunda{} \\
Norte & 245 & 3 horas \referenciaSegunda{} \\
Noroeste & 160 & 3 horas \referenciaSegunda{} \\
Oriental & 180 & 3 horas \referenciaSegunda{} \\
Occidental & 20 & 3 horas \referenciaSegunda{} \\
Peninsular & 140 & 3 horas \referenciaSegunda{} \\
Total & 935 & -- \referenciaSegunda{} \\
\bottomrule
\end{longtable}


\subsubsection*{19. Anexos}


La Convocatoria original prevé un apartado de anexos como parte de su estructura formal.\referenciaBase{}


Para fines de esta versión consolidada de consulta, los anexos deben leerse conforme a la Convocatoria original y, en su caso, a las referencias expresas que hace el numeral 17 vigente al Anexo 19.3, Formato Aceptación de Obras.\referenciaBase{}\referenciaSegunda{}


\subsubsection*{Transitorios de las modificaciones}


\paragraph*{Primera modificación}
\mbox{}\par
\nopagebreak[4]


ÚNICO. El Acuerdo por el que se emitió la modificación publicada el 26 de mayo de 2026 entró en vigor el mismo día de su publicación en el Diario Oficial de la Federación.\marcaPrimera{}\referenciaPrimera{}


\paragraph*{Segunda modificación}
\mbox{}\par
\nopagebreak[4]


PRIMERO. El Acuerdo por el que se emitió la segunda modificación publicada el 10 de julio de 2026 entró en vigor el mismo día de su publicación en el Diario Oficial de la Federación.\marcaSegunda{}\referenciaSegunda{}


SEGUNDO. Para el caso de aquellas personas solicitantes que hayan presentado registros de manifestación de interés y solicitudes de estudios de interconexión y conexión previos a la fecha de publicación de ese Acuerdo, en un plazo no mayor a 5 días hábiles, debieron manifestar su interés en continuar con el procedimiento de la Convocatoria de Proyectos Estratégicos conforme al procedimiento establecido previo a la entrada en vigor del propio Acuerdo.\marcaSegunda{}\referenciaSegunda{}


En caso de que las personas solicitantes no manifestaran ese interés, el procedimiento continuaría conforme a lo establecido en la segunda modificación.\marcaSegunda{}\referenciaSegunda{}

